\documentclass[12pt]{article}

\usepackage{graphicx}
\usepackage{amsmath}
\usepackage[colorlinks]{hyperref}
\usepackage{endfloat}

\title{Guía 1}
\author{Joaquín Rapela}

\begin{document}

\maketitle

\section*{Problema 6}
\subsection*{Parte a}

La expresión del campo eléctrico sobre el eje $x$ es:

\begin{align}
    E(x, 0, 0)&=\frac{q}{4\pi\epsilon_0}\left(\frac{x+d/2}{|x+d/2|^3}-\frac{x-d/2}{|x-d/2|^3}\right)\hat{i}
\end{align}

La figura~\ref{fig:6a1_3} grafica $E(x, 0, 0)$ en función de $x$.  El sentido
del $E(x, 0, 0)$ varía a lo largo del eje $x$. La figura~\ref{fig:6a5} grafica
$||E(x, 0, 0)||_2$ en función de $x$.


\begin{figure}
    \centering
    \href{https://www.gatsby.ucl.ac.uk/~rapela/em/delfi/mySolutions/guia1/figures/6a1_3.html}{\includegraphics[width=6in]{../figures/6a1_3.png}}

    \caption{E(x, 0, 0) para $q=0.1\mu C$, y $d=8 m$. Haga clic en la imagen
    para ver su versión interactiva. El código usado para construir esta figura
    se encuentra
    \href{https://github.com/joacorapela/em/blob/master/delfi/mySolutions/guia1/code/scripts/6a1_3.py}{aquí}.}

    \label{fig:6a1_3}
\end{figure}

\begin{figure}
    \centering
    \href{https://www.gatsby.ucl.ac.uk/~rapela/em/delfi/mySolutions/guia1/figures/6a5.html}{\includegraphics[width=6in]{../figures/6a5.png}}

    \caption{$||E(x, 0, 0)||_2$ para $q=0.1\mu C$, y $d=8 m$. Haga clic en la
    imagen para ver su versión interactiva. El código usado para construir esta
    figura se encuentra
    \href{https://github.com/joacorapela/em/blob/master/delfi/mySolutions/guia1/code/scripts/6a5.py}{aquí}.}

    \label{fig:6a5}
\end{figure}

\subsection*{Parte b}

La expresión del campo eléctrico sobre el eje $y$ es:

\begin{align}
    E(0, y, 0)&=\frac{qd}{(4\pi\epsilon_0)((d/2)^2+y^2)^{3/2}}\hat{i}
\end{align}

La figura~\ref{fig:6b} grafica $E(0, y, 0)$ en función de $y$.  El sentido
del $E(0, y, 0)$ no varía a lo largo del eje $y$.


\begin{figure}
    \centering
    \href{https://www.gatsby.ucl.ac.uk/~rapela/em/delfi/mySolutions/guia1/figures/6b.html}{\includegraphics[width=6in]{../figures/6b.png}}

    \caption{E(0, y, 0) para $q=0.1\mu C$, y $d=8 m$. Haga clic en la imagen
    para ver su versión interactiva. El código usado para construir esta figura
    se encuentra
    \href{https://github.com/joacorapela/em/blob/master/delfi/mySolutions/guia1/code/scripts/6b.py}{aquí}.}
    \label{fig:6b}
\end{figure}

\subsection*{Parte c}

La expresión del campo eléctrico sobre el eje $y$ es:

\begin{align*}
    E(\mathbf{p})&=\frac{q}{4\pi\epsilon_0}
    \left(
    \frac{\mathbf{p}-\begin{bmatrix}
                       -d/2\\
                       0\\
                       0
                     \end{bmatrix}}
         {\left\Vert\mathbf{p}-\begin{bmatrix}
                                 -d/2\\
                                 0\\
                                 0
                               \end{bmatrix}\right\Vert^3}-
    \frac{\mathbf{p}-\begin{bmatrix}
                       d/2\\
                       0\\
                       0
                     \end{bmatrix}}
         {\left\Vert\mathbf{p}-\begin{bmatrix}
                                 d/2\\
                                 0\\
                                 0
                               \end{bmatrix}\right\Vert^3}
    \right)
\end{align*}

La figura~\ref{fig:6c} grafica $E(\mathbf{p})$ en el espacio.


\begin{figure}
    \centering
    \href{https://www.gatsby.ucl.ac.uk/~rapela/em/delfi/mySolutions/guia1/figures/6c.html}{\includegraphics[width=6in]{../figures/6c.png}}

    \caption{$E(\mathbf{p})$ para $q=0.1\mu C$, y $d=8 m$. Haga clic en la imagen
    para ver su versión interactiva. El código usado para construir esta figura
    se encuentra
    \href{https://github.com/joacorapela/em/blob/master/delfi/mySolutions/guia1/code/scripts/6c.py}{aquí}.}
    \label{fig:6c}
\end{figure}

\section*{Problema 7}
\subsection*{Parte a}

Le siguiente ecuación representa el campo eléctrico sobre en la posición
$(x_0,y_0,y_0)$ generado por una distribución uniforme de cargas en el
intervalo $(-L/2,L/2)$ del eje $x$, con densidad $\lambda$.

\begin{align*}
    E(x_0,y_0,z_0|L)&=\frac{\lambda}{4\pi\epsilon_0}\left(\frac{1}{\sqrt{(x_0-L/2)^2+y_0^2+z_0^2}}-\frac{1}{\sqrt{(x_0+L/2)^2+y_0^2+z_0^2}}\right.,\\
                                                     &\quad\quad\quad\quad\quad\frac{y_0}{y_0^2+z_0^2}\left(\frac{L/2-x_0}{\sqrt{(L/2-x_0)^2+y_0^2+z_0^2}}+\frac{L/2+x_0}{\sqrt{(L/2+x_0)^2+y_0^2+z_0^2}}\right),\\
                                                     &\quad\quad\quad\quad\quad\left.\frac{z_0}{y_0^2+z_0^2}\left(\frac{L/2-x_0}{\sqrt{(L/2-x_0)^2+y_0^2+z_0^2}}+\frac{L/2+x_0}{\sqrt{(L/2+x_0)^2+y_0^2+z_0^2}}\right)\right)
\end{align*}

\subsection*{Partes b, c}

\begin{align*}
    \lim_{L\rightarrow\infty}E(x_0,y_0,z_0|L)&=\frac{\lambda}{4\pi\epsilon_0}\frac{2}{y_0^2+z_0^2}(0,y_0,z_0)
\end{align*}

La figura~\ref{fig:7bc} grafica $||E(0,y_0,0|L)||_2$ en función de $y_0$,
usando $\lambda=15\mu C/m$. Las barras azules corresponden a $L=0.5m$ y las
barras rojas a $L=\infty$.

\begin{figure}
    \centering
    \href{https://www.gatsby.ucl.ac.uk/~rapela/em/delfi/mySolutions/guia1/figures/7bc.html}{\includegraphics[width=6in]{../figures/7bc.png}}

    \caption{$||E(0,y_0,0)||_2$ en función de $y_0$, usando $\lambda=15\mu
    C/m$ y $L=0.5m$ (barras azules) o $L=\infty$ (barras rojas). Haga clic en
    la imagen para ver su versión interactiva. El código usado para construir
    esta figura se encuentra
    \href{https://github.com/joacorapela/em/blob/master/delfi/mySolutions/guia1/code/scripts/7bcd.py}{aquí}.}

    \label{fig:7bc}
\end{figure}

El campo eléctrico correspondiente a $L=\infty$ es siempre más fuerte que el
correspondiente a $L=0.5m$.

\subsection*{Parte d}

La figura~\ref{fig:7d} grafica las diferencias porcentuales entre $\Vert
E(0,y_0,0|L=0.5m)\Vert_2$ y $\lim_{L\rightarrow\infty}\Vert
E(0,y_0,0|L)\Vert_2$, en función de $y_0$.

\begin{figure}
    \centering
    \href{https://www.gatsby.ucl.ac.uk/~rapela/em/delfi/mySolutions/guia1/figures/7d.html}{\includegraphics[width=6in]{../figures/7d.png}}

    \caption{Diferencias porcentuales entre el uso de $L=0.5m$ y $L=\infty$
    (i.e.,
    $\frac{E_{L=0.5m}(0,y_0,0)-E_{L=\infty}(0,y_0,0)}{E_{L=\infty}(0,y_0,0)}$).
    Haga clic en la imagen para ver su versión interactiva.  El código usado
    para construir esta figura se encuentra
    \href{https://github.com/joacorapela/em/blob/master/delfi/mySolutions/guia1/code/scripts/7bcd.py}{aquí}.}

    \label{fig:7d}
\end{figure}

La diferencia porcentual entre $\Vert E(0,y_0,0|L=0.5m)\Vert_2$ y
$\lim_{L\rightarrow\infty}\Vert E(0,y_0,0|L)\Vert_2$ es mayor a medida que
$y_0$ aumenta, ya que en estos casos la influencia de cargas
correspondientes a $|x|>L/2$ es porcentualmente mayor.

\end{document}
